% =====================================================================================
%  Chapter 5 — Experimental Setup and Protocol
%  Thesis: Comparison of Control Strategies for Interaction-Force Aware Multirotor Teams
%
%  DRAFT, 23 August 2026 (writing-track item W.7).
%  Sources: docs/05_Experimental_Protocol_2Robot.md, docs/10_Formation_Library.md,
%           docs/11_Hardware_Readiness_Checklist.md, docs/09_Simulation.md.
%
%  HONESTY NOTE. At the time of writing NOTHING HAS FLOWN. This chapter describes a design,
%  not a completed campaign. Every statement about what "is" done must be checkable; every
%  statement about what will be done is written as such. Do not let the future tense slip
%  into the past as chapters 6-7 are drafted -- update this chapter instead.
%
%  Citation baseline: 2026-08-27 (see docs/thesis/CITATION_BASELINE.md).
%  After that date: only NEW cites or EDITED prior-work sentences need re-checking.
% =====================================================================================

\chapter{Experimental Setup and Protocol}
\label{ch:setup}

Chapter~\ref{ch:control} described seven strategies and argued that four of them differ only in
which residual term is active. That argument is only worth as much as the experiment that exercises
it. This chapter describes the platform, the scenarios, the measurement chain and the protocol, and
in each case the reasoning is the same: the design exists to remove alternative explanations for an
observed difference.

Section~\ref{sec:setup-validity} states the threats this design is built against, and which of
them it does not eliminate.

% -------------------------------------------------------------------------------------
\section{Platform}
\label{sec:setup-platform}

The vehicles are Crazyflie~2.1 brushless quadrotors, nominally \SI{41}{\gram} and
\SI{9}{\centi\metre} rotor-to-rotor. All vehicles in a team are the same platform: the gain block
is shared across robots, so a mixed team would fly at least one vehicle on the wrong mass, thrust
constant and inertia --- which would appear in the logs as a residual and be indistinguishable from
an interaction force (Section~\ref{sec:model-measure-fail}).

Two vehicles are the minimum for the protocol. A third is required for the three-robot scenarios,
which cannot be run without it.

Each vehicle carries two additions that are not optional for this work:

\paragraph{A rotor-speed source} --- optical RPM deck or DShot telemetry. Without it the applied
thrust cannot be reconstructed, so $a_{\mathrm{res}}$ is unavailable and \emph{reads identically
zero} rather than failing loudly. Strategies 1 and 4 additionally cannot run at all. This is the
single most consequential hardware prerequisite in the setup.

\paragraph{A micro SD deck} --- the recording medium for the dataset, for the reasons in
Section~\ref{sec:setup-logging}.

State estimation uses the on-board extended Kalman filter, fed with external pose from a
motion-capture system at \SI{100}{\hertz}. The control loop runs at \SI{500}{\hertz} onboard.

% -------------------------------------------------------------------------------------
\section{Flight volume}
\label{sec:setup-volume}

The tracked volume measures approximately \SI{2}{\metre} $\times$ \SI{4}{\metre} horizontally and
extends to \SI{1.7}{\metre} in height. Two distinct limits are enforced and should not be
conflated: the \emph{tracked volume}, outside which the motion-capture system cannot see the
vehicles, and an \emph{operational floor} of \SI{0.30}{\metre}, chosen so that ground effect does
not contaminate the measurement of inter-vehicle interaction.

The floor is lowered to \SI{0.10}{\metre} for one scenario only --- the near-ground pass of
Section~\ref{sec:setup-scenarios}, whose purpose is to measure ground effect in isolation so that
it can be separated from downwash rather than confounded with it.

The volume is a binding constraint rather than a formality. Several scenarios sit within
approximately \SI{10}{\centi\metre} of a wall at their extremes, and the scenario definitions are
automatically centred and, where necessary, rotated to fit. That a configuration fits is checked
before flight, not discovered during it.

\todo{The volume figures are corroborated from two independent statements but have not been
tape-measured. Measure before the first session and update this section
(docs/11\_Hardware\_Readiness\_Checklist.md).}

% -------------------------------------------------------------------------------------
\section{The formation library}
\label{sec:setup-scenarios}

Sixteen scenarios, fixed before any data collection and not modified afterwards. Freezing them is
what allows every strategy to meet identical excitation; a library that evolved during the campaign
would make later strategies incomparable with earlier ones.

The scenarios are parameterised rather than hard-coded --- separations, offsets and speeds are
arguments --- so that a separation sweep is the same scenario at different parameters rather than
a different scenario. Each is compiled to a per-vehicle polynomial trajectory and uploaded through
the standard trajectory interface, so that the same command produces the same motion in simulation
and on hardware.

\begin{table}[t]
  \centering
  \caption{The formation library. Scenarios marked $\dagger$ carry elevated risk and are gated
    behind an explicit acknowledgement; those marked $\ast$ are null conditions in which no
    inter-vehicle interaction is expected.}
  \label{tab:setup-scenarios}
  \small
  \begin{tabular}{@{}llp{0.46\linewidth}@{}}
    \toprule
    & Scenario & Purpose \\
    \midrule
    \multicolumn{3}{@{}l}{\emph{Two vehicles}}\\
    A1 & Vertical stack, hover      & Baseline. Nothing moves, so any relative displacement is interaction rather than tracking error \\
    A2 & Vertical stack, tracking   & Constant separation while both vehicles move \\
    A3 & Static top                 & Fixed wash source, lower vehicle sweeps through: captures the transition, not only the level \\
    A4 & Offset stack               & Partial overlap; asymmetric and hardest to model \\
    A5 & Reverse-circle stack       & Lateral offset swept twice per lap at constant height and speed \\
    A6 & Extreme stack $\dagger$    & Strongest interaction, smallest error budget \\
    A7 & Dynamic merge $\dagger$    & One continuous sweep of the whole separation range: gives a model the gradient, not only levels \\
    A8 & Vertical swap              & Wash arrives as a transient with a known arrival time; hardest case for a feedforward model \\
    \midrule
    \multicolumn{3}{@{}l}{\emph{Three vehicles}}\\
    B1 & I-stack                    & Lowest vehicle in the combined wash of two \\
    B2 & V-stack                    & Asymmetric aggregation, which B1 cannot produce \\
    B3 & Vertical swap              & Three crossings at three height differences in one flight \\
    \midrule
    \multicolumn{3}{@{}l}{\emph{Controls and coverage}}\\
    C1 & Side-by-side coplanar $\ast$ & A residual here, and not in A1, is not downwash \\
    C2 & Leader--follower line $\ast$ & Wake rather than wash \\
    C3 & Equilateral triangle $\ast$  & Coplanar three-body control \\
    C4 & Docking approach $\dagger$   & Simultaneous lateral and vertical closure \\
    C5 & Near-ground pass, one vehicle & Ground effect alone, separable from downwash \\
    \bottomrule
  \end{tabular}
\end{table}

\subsection{The null scenarios}
\label{sec:setup-null}

C1--C3 are the most important scenarios in the library and the least interesting to watch. They
place the vehicles in configurations where no significant inter-vehicle interaction should exist,
and they exist because of the property established in Section~\ref{sec:model-dynamics}: the
residual is a catch-all, non-zero even in isolated flight, so a residual measured in formation
cannot be attributed to the interaction without a reference.

They serve two purposes. They bound the measurement floor, so that a reported improvement can be
distinguished from noise. And they detect model error: a residual that appears when no neighbour is
in the wash is a mis-identified mass or thrust constant, not an interaction, and must be corrected
before the campaign proceeds rather than explained afterwards.

\subsection{Regimes, and why the library spans them}

The scenarios are not a list of poses but a deliberate spread across the interaction regimes that
Research Question~2 distinguishes. A1 and C1--C3 are static; A2 and A3 are quasi-static; A4 and A5
sweep lateral offset; A7 sweeps separation continuously; A8 and B3 present the disturbance as a
transient with a known arrival time. A model fitted on one regime and evaluated on another is what
makes ``outside the training distribution'' a measurable condition rather than a caveat.

\begin{quote}
\textbf{Requirement for data collection: the set must excite lateral relative motion.}
A purely vertical scenario holds the lateral components of the relative state constant, so a model
trained only on such data has seen no variation in half of its six inputs and is extrapolating the
first time a formation moves sideways. \textbf{The collection set must therefore include A4, A5
and A7 --- not only the vertical stacks A1, A2 and A3.} This is not a preference about coverage; a
set without them produces a model that cannot be evaluated on the scenarios it will meet.
\end{quote}

This was established in simulation before any flight: a dry run collecting only a vertical stack
produced training data in which two of the six network inputs never varied, which the training
pipeline reports as a degenerate normalisation statistic and which no amount of training can
repair.

% -------------------------------------------------------------------------------------
\section{Measurement and logging}
\label{sec:setup-logging}

\subsection{Why the dataset is recorded onboard}

Radio telemetry saturates at approximately two vehicles: each drone transmits on the order of
\num{800} packets per second at full rate against a per-dongle ceiling near \num{1000}. Recording
the dataset over radio would therefore couple the sampling rate to the team size, which is
unacceptable for a study whose subject is what happens as vehicles are added.

The micro SD deck records \num{35} variables at \SI{500}{\hertz} on each vehicle independently of
radio and of team size. Radio logging is retained for live monitoring and for single-vehicle
tuning, where the bandwidth exists.

\subsection{What is recorded}

Position and velocity estimates; attitude; gyro and accelerometer; the measured residual
$a_{\mathrm{res}}$; the INDI torque and filtered angular acceleration; the network prediction
$\hat{a}_{\mathrm{res}}$ and its saturation flag; the commanded setpoint; and the four motor
commands.

Two of these deserve comment. The measured residual is recorded \emph{in every controller mode},
including those that make no use of it, because the training data is collected under the
uncompensated baseline. And the network prediction is recorded \emph{whether or not it is enabled},
because comparing prediction against measurement is how a learned model is evaluated, and that
comparison is only meaningful on a flight where the prediction is not also changing the vehicle's
behaviour.

The logging configuration is fixed before data collection and shared by the collection and
evaluation flights. A configuration change between the two would silently make the two halves of
the dataset incomparable.

\subsection{Time alignment across vehicles}

Each vehicle timestamps its samples against its own boot, so the logs share no clock. Recording is
started by a single broadcast, which aligns the vehicles to within the broadcast jitter plus one
logging period. Residual error comes from propagation, task wake phase and crystal drift between
vehicles; the last is a rate error and accumulates over a flight.

The merge tool therefore \emph{measures} the offset by cross-correlating each vehicle's onboard
trace against a radio trace sharing a common clock, and reports offset and drift rather than
assuming them. The predicted alignment is a few milliseconds; that figure is to be confirmed
against real logs on the first two-vehicle flight.

% -------------------------------------------------------------------------------------
\section{Protocol}
\label{sec:setup-protocol}

\subsection{Hardware gate}
\label{sec:setup-gate}

No data collected before the following sequence passes is used. The gate is not an experiment; it
is the check that nothing is broken, so that what is measured afterwards means something.

\begin{enumerate}
  \item \textbf{Validate.} Single vehicle, geometric first and then INDI, on the trajectory ladder.
        Confirm the gains read back as configured.
  \item \textbf{Tune, then freeze.} Any remaining attitude tuning happens here and only here. The
        gains are then frozen for the remainder of the campaign.
  \item \textbf{Two vehicles at large separation.} Confirm $a_{\mathrm{res}}$ is non-zero and
        correctly signed; confirm the measured time alignment.
\end{enumerate}

The order carries the reasoning. Tuning before validating means tuning against an unknown fault.
Flying two vehicles before the freeze means the gains move underneath the dataset.

Stage~1 cannot confirm that the compensation terms are \emph{correct}, only that they are not
catastrophic: with no neighbour present $a_{\mathrm{res}} \approx 0$, so a compensation term is
dormant and a clean single-vehicle flight demonstrates nothing about it. Stage~3 is the first point
at which the sign of the compensation is observable at all, which is why it is part of the gate
rather than the first experiment.

\subsection{Data collection}
\label{sec:setup-collect}

Collection flights are flown under the \emph{uncompensated} baseline of
Section~\ref{sec:ctrl-s0}. This is not a convenience: a compensating controller partly cancels the
disturbance, so the residual it records is the one that survived compensation, which is not the
quantity a predictive model should be fitted to.

Separations are approached from large to small, and a tighter configuration is attempted only after
the looser ones are clean.

\textbf{The set must span both vertical and lateral excitation} --- A1/A2/A3 \emph{and} A4/A5/A7
--- for the reason given in Section~\ref{sec:setup-scenarios}. A collection campaign that omits the
laterally excited scenarios cannot support the predictive strategies, and the omission is not
recoverable without re-flying.

\subsection{Comparison campaign}
\label{sec:setup-campaign}

Each strategy flies the frozen library. Because the control law is selected at runtime
(Section~\ref{sec:ctrl-runtime}), strategies are interleaved within a session rather than flown on
separate days, which pairs them against the same battery state, temperature and calibration.
Each condition is repeated to allow a distribution rather than a single number to be reported.

The comparison is tightest among Strategies 0, 1, 2 and 4, which share the entire position and
attitude loop and differ only in which residual term is active; a shared gain set is meaningful for
them in a way it is not across structurally different controllers. Strategies 3 and 6 are compared
under identical trajectories, formations and hardware, but not under an identical control
structure, and the discussion chapter must phrase their results accordingly.

\subsection{Metrics}
\label{sec:setup-metrics}

Three axes, corresponding to the research questions.

\paragraph{Tracking error.} Position RMSE, total and per axis, and the maximum deviation from the
commanded trajectory.

\paragraph{Residual rejection.} The error in the \emph{separation between vehicles}, measured
against the commanded separation. Downwash acts as a bias --- the lower vehicle sags, closing the
gap --- so the mean error is reported as the primary quantity and the spread as secondary. A
compensator that works moves the mean back toward the commanded value. Reporting only the spread
would hide precisely the effect being studied.

Alongside it, the magnitude of the measured residual that remains after compensation, and for every
strategy carrying a learned model, the prediction against the measurement.

\paragraph{Robustness and cost.} Control effort; behaviour as separation tightens; abort and
failure rate; onboard computation per control step; sensing prerequisites; and the quantity of
training data required to reach a given prediction accuracy.

% -------------------------------------------------------------------------------------
\section{The role of simulation}
\label{sec:setup-sim}

A software-in-the-loop simulator runs the same controller source that flies, with a published
downwash model providing the interaction force. It was used throughout development, and its role in
this thesis is deliberately limited.

\textbf{Simulation is used to disprove implementations, never to validate a method.} It established
that every scenario in the library executes and produces the commanded geometry under both
controllers, and it exposed two defects in flight code that single-vehicle testing could not have
found. It cannot support a claim about how well a method works, for a reason specific to this
subject: the interaction force in simulation is produced by a learned model of the same family as
the one being fitted, and learning a network's output with another network is a substantially
easier problem than learning real air.

No result in the two-robot results chapter or~the multi-robot results chapter rests on simulation.

% -------------------------------------------------------------------------------------
\section{Threats to validity}
\label{sec:setup-validity}

\begin{table}[t]
  \centering
  \caption{Threats to validity and the design decisions taken against them.}
  \label{tab:setup-threats}
  \small
  \begin{tabular}{@{}p{0.28\linewidth}p{0.34\linewidth}p{0.30\linewidth}@{}}
    \toprule
    Threat & Mitigation & Residual risk \\
    \midrule
    Tuning confound & Gains frozen before collection; shared across strategies; runtime selection allows interleaving & Shared gains may suit some strategies better. Justified for 1/2/4, which share the whole loop; weaker for 3 and 6, which do not (§\ref{sec:ctrl-fairness}) \\
    Measurement floor & Null scenarios C1--C3 bound it under the same pipeline & --- \\
    Model error read as interaction & Null scenarios detect it; airframe constants identified once and frozen & Battery discharge within a flight is not compensated \\
    Training/test leakage & Held-out scenarios; split by contiguous flight segment, never by sample & --- \\
    Session effects & Strategies interleaved within a session; repeated conditions & --- \\
    Silent loss of the residual & $a_{\mathrm{res}} \neq 0$ verified at the start of every session & --- \\
    Simulation-to-hardware gap & No claim rests on simulation & --- \\
    Single airframe & --- & \textbf{Not mitigated.} Generalisation beyond this platform is discussed, not demonstrated \\
    \bottomrule
  \end{tabular}
\end{table}

Splitting the training and test data by contiguous flight segment rather than by sample deserves a
note, because the obvious alternative is badly wrong here. At \SI{500}{\hertz} consecutive samples
are nearly identical, so a random per-sample split places a near-duplicate of almost every test
point in the training set and reports a validation error that means nothing.

Two threats are not mitigated and are stated as limitations rather than managed. The results
concern a single airframe; quantities are normalised where a normalisation is defensible, but
generalisation to other platforms is argued rather than shown. And battery discharge changes the
thrust constant slightly within a flight, which appears as a slow drift in the measured residual;
the effect is bounded by the null scenarios but not removed.

% -------------------------------------------------------------------------------------
\section{Summary}
\label{sec:setup-summary}

Two or three identical brushless quadrotors, each carrying a rotor-speed source and an SD card,
flying sixteen frozen scenarios inside a tracked volume of roughly \SI{2}{\metre} $\times$
\SI{4}{\metre} $\times$ \SI{1.7}{\metre}, with \num{35} variables recorded onboard at
\SI{500}{\hertz} independently of team size. Gains are frozen before any data is collected;
strategies are selected at runtime and interleaved within sessions; and three null scenarios
establish what a measurement of zero looks like under the same pipeline.

\todo{This chapter describes a design. At the time of writing no flight has taken place. Revise to
past tense as the campaign proceeds, and record deviations from this protocol here rather than only
in the results.}
